\paragraph{A mass-twelve gluing lemma.}
Throughout, $D=q+2+q^{-1}$ and $T_j=q^j+q^{-j}$.
Every effective $W\ge D$ with $W\equiv D\pmod2$ and $W(4)=12$
is one of
\[
 D+8,\quad D+4+2T_j,\quad D+4T_j,\quad
 D+2T_i+2T_j\ (i<j).
\]
Indeed $(W-D)/2$ is an effective symmetric character of mass four,
and these are its complete partitions into a central coefficient
and paired nonzero exponents. Hence $W(0)\in\{-8,0,8\}$.
If $W(0)=0$, write $W=DH$. Then
\[
 H\in\mathbb Z[D],\quad H(4)=3,\quad H\equiv1\pmod2,
 \quad H(0)\ne1,-1.
\]
Either $H=3$ (the case $W=3D$), or $H$ is primitive and
irreducible over $\mathbb Q[D]$, with leading coefficient two.
For the last assertion, every nonconstant case has leading
coefficient two and odd constant coefficient. In a nontrivial
primitive integral factorization, one factor has odd leading
coefficient and would keep its positive degree modulo two,
contradicting that the product reduces to one.
For the assertion at zero, the zero-valued cases in the list are
$D+4+2T_o$ with $o$ odd and $D+2T_o+2T_e$ with $o$ odd and
$e$ even. The derivative identity from (G036.R5) gives
\[
 H(0)=1+2o^2\quad\hbox{or}\quad 1+2(o^2-e^2).
\]
Neither is one or minus one: equality would require equal
squares of opposite parity or $e^2-o^2=1$, impossible for
positive distinct integers $e,o$. This also includes $o=1$
and the constant $H=3$.

\paragraph{The nonzero-value branch.}
For any effective $U\ge D$ with $U\equiv D\pmod2$ and mass $u$,
\[
 |U(0)|\le u-4,
\]
because $(U-D)/2$ is effective and its absolute value at $q=-1$
is at most its mass. In every remaining mass-twelve row of
NN-G036, $(u_1-4)(u_2-4)<64$, as follows directly from the
complete listed mass tuples. If $W(0)=8$ or $-8$, the identities
\[
 WG_1-U_1G_2=DW\delta_1,\qquad
 WG_2-U_2G_1=DW\delta_2
\]
have an invertible coefficient matrix at $D=0$ over $\mathbb Q$.
Thus $D$ divides both $G_i$ in $\mathbb Z[D]$.
Using $G_1\equiv0$ and $G_2\equiv D^2(D-1)\pmod2$, their
quotients by $D$ have even evaluation at $D=4$. Consequently
both $g_i$ are divisible by eight. This excludes any mass row
containing $84$ in this branch.
If also $\delta_1=\delta_2=\psi=D(D-3)$, divide the same
identities by $D$ and apply the invertible matrix at zero once
again. Both $G_i$ are divisible by $D^2$. Since $WS=G_1G_2$
and $W(0)\ne0$, the polynomial $S$ is divisible by $D^4$.
In particular, the mass block $(g_1,g_2)=(80,96)$ or $(96,80)$,
where $s=640$, is impossible: $4^4$ does not divide $640$.

\paragraph{The two blocks with action masses $84,88$.}
Their action masses are $84,88$ in either order, with associated
$u$-masses six and eight. The preceding paragraph leaves only
$W=DH$. If $H$ is nonconstant, it cannot divide the mass-six
character $D+2$ (the integral leading coefficient would have
to divide one), nor a mass-eight character (its value three
would have to divide eight). The common-factor argument of
(G036.R2) forces $H$ to divide both $G_i$, contradicting that
three does not divide $88$.
If $H=3$, reduction modulo three supplies the same argument
in the integral domain $\mathbb F_3[D]$. Both $U_i$ have
nonzero reduction: one is $D+2$, and the other has mass eight.
The three identities (G036.R2) then force both $G_i$ to have
all coefficients divisible by three, again contradicting
mass $88$. Thus both base rows are impossible.

\paragraph{The block with both action masses $84$.}
Here $U_1=U_2=D$ and $d_1=d_2=14$. Again only $W=DH$
remains. The common-factor argument (or reduction modulo
three when $H=3$) gives $G_i=HP_i$ integrally. Their masses
are $P_i(4)=28$, and their reductions are
$P_1\equiv0$, $P_2\equiv D^2(D-1)\pmod2$.
Cancelling $DH$ in the two cell identities gives
\[
 P_2=HP_1-D\delta_1,\qquad
 P_1=HP_2-D\delta_2.
\]
Since $H(0)^2\ne1$, evaluating at zero forces
$P_1(0)=P_2(0)=0$. In particular $P_1/D$ is an integral
polynomial with every coefficient even, but its evaluation
at four is $28/4=7$. This is a contradiction.

\paragraph{The two smallest mass-twelve rows.}
Suppose $(g_1,g_2)$ is $(80,96)$ or $(96,80)$ and both
cells are $\psi$. Let $G_-$ have mass eighty and $G_+$
have mass ninety-six; their $U$-masses are eight and twelve.
The nonzero-value branch was already excluded, so $W=DH$.
For nonconstant $H$, its value three cannot divide eighty,
so $H\nmid G_-$. The identities (G036.R2) force
$H\mid G_+$ and $H\mid U_+$. Every mass-twelve $U_+$
has value at zero equal to zero or $\pm8$. Since $H(0)$ is
odd and is neither one nor minus one, its divisibility
forces $U_+(0)=0$. Write $U_+=DJ$. The classification above
applies to $J$: it is either three or primitive irreducible
of leading coefficient two. Since $H\mid J$, the latter
case is forced and $J=H$. Thus $U_+=W$.
For $H=3$, reduce $U_+G_-=3D(G_+-D\psi)$ modulo three.
Since $G_-$ has mass eighty, its reduction is nonzero, so
all coefficients of $U_+$ are divisible by three. Its odd
coefficients at $q,q^{-1}$ are at least three each, and its
even central coefficient is at least six. These exhaust its
mass twelve, so $U_+=3D=W$ also in this case.

The larger-cell identity now gives $G_+=G_-+D\psi$.
The other identity becomes
\[
             (W-U_-)G_-=D\psi(W+U_-).
\]
The mass-eight classification gives $U_-(0)=4$ or $-4$,
while $W(0)=0$. Hence $D^2\mid G_-$.
If $G_-=G_1$, its quotient by $D^2$ has even coefficients
but mass $80/16=5$, a contradiction.
If $G_-=G_2$, use the necessary mod-four identity
$G_2-D\delta_2\equiv G_1\pmod4$ from NN-G030.
Here $\delta_2=\psi$ and $G_1=G_2+D\psi$, so it would
require $2D\psi\equiv0\pmod4$, which is false.
This excludes both rows, including every auxiliary increment.

These five exclusions use neither a degree bound nor
positivity of any quotient. They retain precisely the
closed-table character dictionary and scope of NN-G036.
